% IEEE/ACM Transactions on Computational Biology and Bioinformatics (TCBB)
% IEEE Computer Society journal template (IEEEtran v1.8b, compsoc mode).
%
% Numbers in the text come from generated/data_numbers.tex and tables/*.tex,
% written by scripts/ from committed results/. Do not type results by hand.
\documentclass[10pt,journal,compsoc]{IEEEtran}

\ifCLASSOPTIONcompsoc
  \usepackage[nocompress]{cite}
\else
  \usepackage{cite}
\fi
\usepackage{amsmath,amssymb}
\usepackage{graphicx}
\usepackage{booktabs}
\usepackage{url}
\usepackage[hidelinks]{hyperref}

\hyphenation{op-tical net-works semi-conduc-tor}

% AUTO-GENERATED by scripts/05_section3_numbers.py -- do not edit by hand.
\newcommand{\nLibraries}{22}
\newcommand{\nGenes}{33,694}
\newcommand{\nBarcodes}{215,960}
\newcommand{\nPublishedCells}{173,564}
\newcommand{\nCloneTableCells}{115,046}
\newcommand{\pctMultiBarcode}{49}
\newcommand{\pctMultiBarcodeMin}{36}
\newcommand{\pctMultiBarcodeMax}{53}
\newcommand{\nStageIVCells}{5,329}
\newcommand{\nStageIVCollisions}{20}
\newcommand{\nExpThreeClones}{555}
\newcommand{\nExpThreeCells}{2,692}
\newcommand{\nExpThreePubClones}{559}
\newcommand{\nExpThreePubCells}{2,816}
\newcommand{\pctExpThreeCloneGap}{0.7}
\newcommand{\pctExpThreeCellGap}{4.4}
\newcommand{\swapRepTwoGSM}{GSM6889993, GSM6889994}
\newcommand{\swapRepThreeGSM}{GSM6889987, GSM6889988}
\newcommand{\nPaperQCPass}{178,226}
\newcommand{\nCloneTableFailPaperQC}{11,350}
\newcommand{\pctCloneTableFailPaperQC}{9.9}
\newcommand{\nCloneTableOverUMICap}{10,973}
\newcommand{\pctStageIVOverUMICap}{40}
\newcommand{\nExpThreeClonesPaperQC}{437}
\newcommand{\umiCap}{30,000}
\newcommand{\nPrimaryCells}{198,073}
\newcommand{\nCloneTablePrimary}{114,667}
\newcommand{\nCloneTableDropped}{379}
\newcommand{\pctCloneTableKept}{99.7}
\newcommand{\nScrubletAutoTotal}{5,896}
\newcommand{\nScrubletAutoMax}{2,434}
\newcommand{\nLibsScrubletZero}{12}
\newcommand{\rhoDoubletUMI}{0.18}
\newcommand{\umiRatioFlagMin}{1.7}
\newcommand{\umiRatioFlagMax}{2.9}
\newcommand{\nSensitivityCells}{181,018}
\newcommand{\nCloneTableSensitivity}{103,049}
\newcommand{\nDACellsCrossStage}{8}
\newcommand{\minGenes}{1,000}
\newcommand{\maxMt}{15}
\newcommand{\dblRatePct}{0.8}
\newcommand{\dblRatePer}{1,000}
\newcommand{\dayOne}{13}
\newcommand{\dayTwo}{21}
\newcommand{\dayThree}{30}
\newcommand{\dayFour}{45}


\begin{document}

\title{Lineage-Grounded Evaluation of Optimal-Transport Couplings and Marker States
in Human Neural Differentiation}

\author{Victor~Ikenna~Kanu
        and~Dong-Seong~Kim%
\IEEEcompsocitemizethanks{\IEEEcompsocthanksitem V.~I.~Kanu and D.-S.~Kim are with the
Network Systems Lab, Kumoh National Institute of Technology, Gumi, South Korea.\protect\\
E-mail: vkanu@kumoh.ac.kr}}

\markboth{IEEE/ACM Transactions on Computational Biology and Bioinformatics}%
{Kanu and Kim: Lineage-Grounded Evaluation of Couplings and Marker States}

\IEEEtitleabstractindextext{%
\begin{abstract}
% Written last, from the final results.
\end{abstract}

\begin{IEEEkeywords}
Single-cell RNA sequencing, lineage tracing, optimal transport, cell-state
forecasting, benchmarking, neural differentiation.
\end{IEEEkeywords}}

\maketitle
\IEEEdisplaynontitleabstractindextext
\IEEEpeerreviewmaketitle

% ---------------------------------------------------------------------------
\IEEEraisesectionheading{\section{Introduction}\label{sec:intro}}
% Two assumptions under test: (1) inferred couplings link cells to their true
% descendants; (2) marker states such as (P, D) capture fate. RQ1--RQ4.

% ---------------------------------------------------------------------------
\section{Related Work}\label{sec:related}

% ---------------------------------------------------------------------------
\section{Data and Setup}\label{sec:data}

\subsection{Lineage-traced differentiation data}
We use the single-cell split-barcoding (SISBAR) data of You \emph{et al.}~\cite{You_2023}:
H9 human embryonic stem cells differentiated towards ventral midbrain and hindbrain
fates and profiled by droplet single-cell RNA sequencing at stages I, II, III and IV
(days~\dayOne, \dayTwo, \dayThree{} and~\dayFour). Each of three experiments labels
progenitors with a retroviral barcode library at one stage, lets them divide three to
four times, and splits the culture so that one half is profiled at that stage and the
other at the next. Experiments~1--3 thus link I$\to$II, II$\to$III and III$\to$IV. A
barcode observed in both halves links \emph{sister} cells rather than one cell observed
twice; this shared ancestry is our ground truth. Because each experiment is barcoded once
and sampled at two consecutive stages, no clone spans three stages, and every analysis
concerns a single transition $t\to t+1$.

Each experiment has three biological replicates, deposited in GEO as GSE207921
(replicate~1) and GSE221592 (replicates~2--3): \nLibraries{} 10x libraries with
\nBarcodes{} cell barcodes and \nGenes{} genes. The authors' barcode-to-cell table lists
\nCloneTableCells{} cells with their viral barcodes and cluster labels. Every cell in it
was found in the 10x library predicted by its barcode suffix and replicate. Stage-IV
barcodes lack the library suffix, so we take the replicate from the viral barcode, which
also resolves \nStageIVCollisions{} stage-IV cell barcodes shared between libraries. In
experiment~3 the GEO sample names swap replicates~2 and~3 relative to the viral barcodes,
consistently in both stages (viral replicate~2: \swapRepTwoGSM; viral replicate~3:
\swapRepThreeGSM). This leaves clone pairing unchanged; we use the viral-barcode replicate
throughout.

\subsection{Clone definition}
Of the clone-table cells, \pctMultiBarcode\% carry more than one viral barcode
(\pctMultiBarcodeMin--\pctMultiBarcodeMax\% across stage--experiment sheets). We define a
clone as the exact set of viral barcodes on a cell, within one experiment and replicate,
and call it cross-stage if it has cells at both stages of its transition. Treating each
barcode as a clone, or merging barcodes that co-occur in a cell, does not reproduce the
published counts; the barcode-set definition reproduces I$\to$II and II$\to$III exactly
(Table~\ref{tab:data_counts}). For III$\to$IV we obtain \nExpThreeClones{} clones and
\nExpThreeCells{} cells against \nExpThreePubClones{} and \nExpThreePubCells{} published
(\pctExpThreeCloneGap\% and \pctExpThreeCellGap\% fewer). All \nStageIVCells{} stage-IV
clone-table cells are present in their expression matrices, so the shortfall lies in the
deposited table. We fixed the handling rule beforehand (one matching pass, then accept the
deposited count) and did not adjust filters to recover the published value.

\subsection{Cell filtering and doublets}
The original study reports thresholds of at least \minGenes{} genes, at most \umiCap{}
UMIs and at most \maxMt\% mitochondrial reads. These pass \nPaperQCPass{} barcodes but
reject \nCloneTableFailPaperQC{} (\pctCloneTableFailPaperQC\%) of the authors' own
clone-table cells, \nCloneTableOverUMICap{} of them on the UMI cap alone, including
\pctStageIVOverUMICap\% of stage-IV clone-table cells; maturing neurons have larger
transcriptomes. Under these thresholds III$\to$IV would retain \nExpThreeClonesPaperQC{}
cross-stage clones (Table~\ref{tab:data_counts}), so the lineage table was evidently not
filtered on UMIs. Because couplings are fitted on barcoded and unbarcoded cells together,
one rule must cover both. Our primary rule keeps cells with at least \minGenes{} genes and
at most \maxMt\% mitochondrial reads, with no UMI cap: \nPrimaryCells{} cells, including
\nCloneTablePrimary{} of \nCloneTableCells{} clone-table cells (\pctCloneTableKept\%).
Neither rule reproduces the published \nPublishedCells{} cells, which reflect further
filtering in the original pipeline.

Doublet detection with Scrublet~\cite{Wolock_2019} was pre-specified per library at an
expected rate of \dblRatePct\% per \dblRatePer{} barcodes, but its automatic threshold was
unstable: it flagged between none and \nScrubletAutoMax{} cells per library (none in
\nLibsScrubletZero{} of \nLibraries{}; \nScrubletAutoTotal{} in total), its output
depended on implementation settings, and observed scores overlapped simulated doublets, as
expected when cells within a stage form a continuum. We therefore do not filter on it.
Instead, each library's top expected-rate fraction by doublet score is flagged. Flagged
cells have \umiRatioFlagMin--\umiRatioFlagMax{} times the median UMI count of other cells,
yet the score correlates only weakly with UMI count (Spearman $\rho=\rhoDoubletUMI$).
Every RQ1--RQ3 analysis is repeated without flagged cells (\nSensitivityCells{} cells;
\nCloneTableSensitivity{} clone-table cells) as a pre-specified sensitivity analysis.

All analyses use one-step pairs from the primary set (Table~\ref{tab:data_counts}).
Dopaminergic neurons are rare among lineage-traced cells: only \nDACellsCrossStage{}
stage-IV dopaminergic cells fall in cross-stage III$\to$IV clones. We therefore study
dopaminergic fate as progenitor fate bias at stage~III rather than as a forecasting target.

% AUTO-GENERATED by scripts/05_section3_numbers.py -- do not edit by hand.
\begin{table}[!t]
\centering
\caption{Cross-stage clones / cells in those clones, per transition and cell set.}
\label{tab:data_counts}
\footnotesize
\setlength{\tabcolsep}{3pt}
\begin{tabular}{lrrr}
\hline
Cell set & I$\to$II & II$\to$III & III$\to$IV \\
 & d13$\to$d21 & d21$\to$d30 & d30$\to$d45 \\
\hline
Published~\cite{You_2023} & 768 / 2,762 & 1,213 / 5,856 & 559 / 2,816 \\
All clone-table cells & 768 / 2,762 & 1,213 / 5,856 & 555 / 2,692 \\
Paper thresholds (not used) & 706 / 2,462 & 1,099 / 5,010 & 437 / 1,989 \\
Primary (this study) & 756 / 2,696 & 1,208 / 5,817 & 549 / 2,667 \\
Doublet sensitivity & 689 / 2,364 & 1,172 / 5,532 & 519 / 2,459 \\
\hline
\end{tabular}
\end{table}


% ---------------------------------------------------------------------------
\section{Methods}\label{sec:methods}
% Protocol (Table 1): clone-level splits, training-only normalisation, 20 seeds,
% paired bootstrap CIs, Wilcoxon--Holm, Cohen's d_z, pre-fixed equivalence margin.
% OT settings (UPDATES.md Q4). Models: Ridge, matched random forest, Transformer.
% Controls: TCBB simulation (positive); independent pairing, within-stage
% clone-label shuffle (negative).

% ---------------------------------------------------------------------------
\section{Results}\label{sec:results}

\subsection{Coupling recovery}\label{sec:rq1}
% RQ1 (Fig. 2).

\subsection{Fate-level agreement}\label{sec:fate}
% RQ1, stage-III progenitors (Fig. 3).

\subsection{Forecasting on true lineage}\label{sec:rq2}
% RQ2 (Fig. 4, Table 2).

\subsection{State adequacy}\label{sec:rq3}
% RQ3 and Jerber transfer (Fig. 5).

\subsection{Consequence for perturbation models}\label{sec:rq4}
% RQ4 (Fig. 6).

% ---------------------------------------------------------------------------
\section{Discussion}\label{sec:discussion}
% Table 3 checklist.

% ---------------------------------------------------------------------------
\section{Limitations and Conclusion}\label{sec:conclusion}

% ---------------------------------------------------------------------------
\ifCLASSOPTIONcompsoc
  \section*{Acknowledgments}
\else
  \section*{Acknowledgment}
\fi
% Funding.

\bibliographystyle{IEEEtran}
\bibliography{sisbar_refs}

\end{document}
